\chapter{La précision, et sa loi impitoyable}

Ce chapitre contient la mauvaise nouvelle de la méthode, puis la raison pour
laquelle on l'emploie quand même — et cette raison est écrasante.

\section{La loi du $1/\sqrt{N}$}

On aimerait que doubler le nombre d'échantillons divise l'erreur par deux. Ce
n'est pas le cas. L'erreur type de l'estimateur vaut :

\[ \sigma_N = \frac{\sigma}{\sqrt{N}} \]

où $\sigma$ est l'écart-type de la fonction échantillonnée. La décroissance se
fait en \textbf{racine carrée} du nombre d'échantillons.

\begin{nkretenir}[Le chiffre à retenir, celui qui explique tout]
\textbf{Pour diviser le bruit par 2, il faut 4 fois plus d'échantillons. Pour
le diviser par 10, il en faut 100 fois plus.}

C'est exactement pourquoi une image en \emph{path tracing} met des heures à
devenir propre : passer d'une image granuleuse à une image lisse ne coûte pas
un peu plus cher, il coûte cent fois plus cher.
\end{nkretenir}

\begin{center}
\begin{tabular}{@{}rll@{}}
\toprule
\textbf{Échantillons} & \textbf{Bruit relatif} & \textbf{Temps (base 1)} \\
\midrule
100        & 10~\%   & 1 \\
400        & 5~\%    & 4 \\
10\,000    & 1~\%    & 100 \\
1\,000\,000& 0,1~\%  & 10\,000 \\
\bottomrule
\end{tabular}
\end{center}

Deux conséquences pratiques, immédiates :

\begin{itemize}
  \item \textbf{inutile d'espérer la perfection par la force}. Multiplier la
        machine par deux ne réduit le bruit que de 30 \%. Le levier n'est pas
        la puissance, c'est la réduction de la variance $\sigma$ (chapitre 4) ;
  \item \textbf{le rendement décroît}. Les premières secondes d'un rendu
        apportent l'essentiel ; les dernières heures polissent. Un moteur qui
        affiche l'image \emph{pendant} qu'elle converge est donc utilisable
        bien avant d'être terminé — c'est le principe du rendu progressif.
\end{itemize}

\section{La bonne nouvelle : la dimension ne compte pas}

Voici pourquoi Monte-Carlo est irremplaçable, et non pas seulement commode.

Les méthodes classiques d'intégration (rectangles, trapèzes, quadrature de
Gauss) découpent le domaine en une grille. En dimension $d$, une grille de $n$
points par axe demande $n^d$ points au total. Regardez la catastrophe :

\begin{center}
\begin{tabular}{@{}rlr@{}}
\toprule
\textbf{Dimension} & \textbf{Points d'une grille $10\times$} & \textbf{Faisable ?} \\
\midrule
1  & 10                  & oui \\
2  & 100                 & oui \\
5  & 100\,000            & oui \\
10 & 10\,000\,000\,000   & non \\
20 & $10^{20}$           & jamais \\
\bottomrule
\end{tabular}
\end{center}

C'est ce qu'on appelle la \emph{malédiction de la dimension}. Et voici le point
décisif :

\begin{nkretenir}[Le $1/\sqrt{N}$ ne dépend pas de la dimension]
L'erreur de Monte-Carlo vaut $\sigma/\sqrt{N}$ \textbf{en dimension 1 comme en
dimension 100}. La dimension n'apparaît nulle part dans la formule.

En rendu, chaque rebond de lumière ajoute deux dimensions (la direction du
rebond). Un chemin à cinq rebonds est une intégrale en dimension 10, et
personne ne le remarque : le code est le même, le coût est le même. Aucune
autre méthode ne tient à ce régime.
\end{nkretenir}

Autrement dit, on ne choisit pas Monte-Carlo parce qu'il est bon — il est lent.
On le choisit parce qu'au-delà de quelques dimensions, \textbf{il est le seul
qui existe encore}.

\section{Estimer sa propre erreur}

Un estimateur honnête sait dire à quel point il est sûr de lui. Comme
l'écart-type $\sigma$ est inconnu, on l'estime à partir des mêmes échantillons :

\[ s^2 = \frac{1}{N-1}\sum_{i=1}^{N}\left(y_i - \bar{y}\right)^2
   \qquad\text{avec}\quad y_i = \frac{f(x_i)}{p(x_i)} \]

et l'erreur type de la moyenne vaut alors $s/\sqrt{N}$.

\begin{nkcode}{Accumulateur qui rend la moyenne ET sa barre d'erreur}
// Algorithme de Welford : une seule passe, et surtout STABLE. La formule
// naive (somme des carres moins carre de la somme) soustrait deux grands
// nombres presque egaux et perd toute la precision -- elle rend meme des
// variances NEGATIVES sur de longues series.
class NkAccumulateurMC {
    public:
        void Ajouter(float64 y) {
            mN += 1;
            const float64 delta = y - mMoyenne;
            mMoyenne += delta / (float64)mN;
            mM2 += delta * (y - mMoyenne);   // le second delta est le NOUVEAU
        }

        float64 Moyenne() const { return mMoyenne; }

        // Erreur type de la moyenne : l'intervalle « moyenne +/- 2 sigma »
        // contient la vraie valeur environ 95 fois sur 100.
        float64 ErreurType() const {
            if (mN < 2)
                return 0.0;
            const float64 variance = mM2 / (float64)(mN - 1);
            return NkMath::Sqrt(variance / (float64)mN);
        }

    private:
        nk_uint64 mN = 0;
        float64 mMoyenne = 0.0;
        float64 mM2 = 0.0;
};
\end{nkcode}

\begin{nkretenir}[À quoi sert la barre d'erreur, concrètement]
Elle permet d'\textbf{arrêter au bon moment}. Plutôt que de fixer un nombre
d'échantillons au hasard, on continue jusqu'à ce que l'erreur type passe sous
un seuil. Un moteur de rendu peut ainsi donner plus d'échantillons aux pixels
bruyants (les zones d'ombre douce, les caustiques) et moins aux pixels déjà
propres — à budget égal, l'image entière converge plus vite.
\end{nkretenir}

\begin{nkpiege}[La barre d'erreur ne voit pas ce qu'elle n'a pas tiré]
Si une contribution énorme se cache dans une région minuscule que vos
échantillons n'ont jamais visitée, votre estimateur sera \emph{confiant et
faux} : peu de variance observée, donc petite barre d'erreur, et pourtant à
côté de la plaque.

C'est le cas classique d'une petite source de lumière très intense. La barre
d'erreur mesure la dispersion de ce que vous avez vu — jamais l'existence de ce
que vous n'avez pas vu.
\end{nkpiege}

\begin{nkexo}[Constater la loi de la racine]
Estimez $\pi$ 50 fois avec $N = 100$, puis 50 fois avec $N = 400$, puis 50 fois
avec $N = 1600$. Calculez à chaque fois l'écart-type de vos 50 estimations.

Vous devez trouver un rapport d'environ 2 entre chaque palier — pas 4. Si vous
trouvez 4, votre générateur aléatoire produit des tirages corrélés (chapitre 7).
\end{nkexo}
